% !TEX program = xelatex
\documentclass[11pt,a4paper,fontset=fandol]{ctexart}
\usepackage[a4paper,left=20mm,right=20mm,top=20mm,bottom=18mm,headheight=15pt]{geometry}
\usepackage{amsmath,amssymb,fancyhdr,lastpage,xcolor}
\definecolor{accent}{HTML}{2C4A7C}\definecolor{problemred}{HTML}{B42318}\definecolor{fixgreen}{HTML}{137333}
\setlength{\parindent}{0pt}\setlength{\parskip}{0.45em}\linespread{1.12}
\pagestyle{fancy}\fancyhf{}\fancyhead[L]{\small 数学二}\fancyhead[R]{\small 第4讲课后题}\fancyfoot[C]{\small 第 \thepage\ 页，共 \pageref{LastPage} 页}\renewcommand{\headrulewidth}{0.6pt}
\newcommand{\pt}[1]{\par\vspace{.3em}{\large\bfseries\color{accent}#1}\par\vspace{.1em}}
\newcommand{\problem}[1]{\textcolor{problemred}{\textbf{【问题】#1}}}
\newcommand{\fix}[1]{\textcolor{fixgreen}{\textbf{【修正】#1}}}
\begin{document}
\section*{4.1}\pt{原题}设函数 $g(x)$ 可微，$h(x)=e^{1+g(x)}$，$h'(1)=1$，$g'(1)=2$，则 $g(1)=(\quad)$。\pt{我的解答}
$g(x)$可微，$h(x)=e^{1+g(x)}$，$h'(1)=1$，$g'(1)=2$，求$g(1)$。\\ $h'(x)=e^{1+g(x)}\cdot g'(x)$\\ 又有$h'(1)=1，g'(1)=2$\\ $\Rightarrow h'(1)=e^{1+g(1)}\cdot2=1$\\ $\Rightarrow g(1)=\ln\dfrac12-1=-\ln2-1$\\ $\Rightarrow C$
\section*{4.2}\pt{原题}设 $f(x)=\displaystyle\lim_{t\to0}x(1-2t)^{-x/t}$，则 $f'(x)=$\underline{\hspace{2cm}}。\pt{我的解答}
$f(x)=\lim_{t\to0}x(1-2t)^{-x/t}=x\cdot\lim_{t\to0}(1-2t)^{-x/t}$\\ $=x\cdot\lim_{t\to0}\exp\left\{-\dfrac{x}{t}\ln(1-2t)\right\}=x\cdot e^{2x}$\\ $\Rightarrow f'(x)=e^{2x}+x\cdot e^{2x}\cdot2=e^{2x}(2x+1)$
\section*{4.3}\pt{原题}已知 $f'(x)=Ae^x$（$A$为正常数），则 $f(x)$ 的反函数的二阶导数为\underline{\hspace{2cm}}。\pt{我的解答}
$f'(x)=Ae^x$，求$f(x)$反函数的二阶导数\\ 二阶导数，再对$f(x)$求导。\quad\problem{题目要求反函数的二阶导数；此处改成了对原函数继续求导。}\\ $\Rightarrow f''(x)=Ae^x\cdot f'(x)=A^2e^{2x}$。\quad\problem{即使求原函数二阶导，也应为 $f''(x)=Ae^x$；此处多乘了 $f'(x)$。}\pt{正确解答}
\fix{设反函数为 $x=x(y)$：}$\displaystyle\frac{dx}{dy}=\frac1{f'(x)}$，$\displaystyle\frac{d^2x}{dy^2}=-\frac{f''(x)}{[f'(x)]^3}=-\frac{Ae^x}{(Ae^x)^3}=-\frac1{A^2e^{2x}}$。
\section*{4.4}\pt{原题}设 $y=y(x)$ 由 $x=\ln(1+t^2)+1$，$y=2\arctan t-(t+1)^2$ 确定，求 $d^2y/dx^2$。\pt{我的解答}
$\dfrac{dx}{dt}=\dfrac{2t}{1+t^2}$，$\dfrac{d^2x}{dt^2}=\dfrac{2(1-t^2)}{(1+t^2)^2}$\\ $\dfrac{dy}{dt}=\dfrac2{1+t^2}-2(t-1)$。\quad\problem{$-(t+1)^2$ 的导数应为 $-2(t+1)$；这是最早错误。}\\ $\dfrac{d^2y}{dt^2}=-\dfrac{4t}{(1+t^2)^2}-2$\\ $\dfrac{d^2y}{dx^2}=\dfrac{\dfrac{d^2y}{dt^2}\dfrac{dx}{dt}-\dfrac{dy}{dt}\dfrac{d^2x}{dt^2}}{(\dfrac{dx}{dt})^3}$，随后代入、通分并得到\\ $=-t^2+\dfrac12t-1-\dfrac1{2t}-\dfrac1{t^3}$。\quad\problem{该结果受上面的导数符号错误影响。}\pt{正确解答}
\fix{只改正 $dy/dt$：}$\displaystyle\frac{dy}{dt}=\frac2{1+t^2}-2(t+1)$，故 $\displaystyle\frac{dy}{dx}=-(t^2+t+1)$，$\displaystyle\frac{d^2y}{dx^2}=-\frac{(1+2t)(1+t^2)}{2t}$。
\section*{4.5}\pt{原题}设 $f'(x)=e^{f(x)}$，$f(2)=1$，求 $f^{(n)}(2)$（$n\ge1$）。\pt{我的解答}
$f'(x)=e^{f(x)}$，$f''(x)=e^{2f(x)}$，$f'''(x)=2e^{3f(x)}$，$f^{(4)}(x)=2\times3e^{4f(x)}$，$\cdots$\\ $\Rightarrow f^{(n)}(x)=(n-1)!e^{nf(x)}$\\ $\Rightarrow f^{(n)}(2)=(n-1)!e^{nf(2)}=(n-1)!e^n$
\section*{4.6}\pt{原题}设 $y=y(x)$ 是由方程 $\sin(xy)=\ln\dfrac{x+e}{y}+1$ 确定的隐函数，求 $y'(0)$。\pt{我的解答}
两边对$x$求导\\ $\Rightarrow\cos(xy)\cdot(y+xy')=\dfrac{y}{x+e}\cdot\dfrac{y-(x+e)y'}{y^2}$\\ $x=0$时，$y=e^2$\\ $\Rightarrow e^2=e\cdot\dfrac{e^2-ey'(0)}{e^4}$\\ $\Rightarrow y'(0)=e-e^4$
\section*{4.7}\pt{原题}已知 $g(x)$ 在 $x=0$ 处二阶可导，且 $g(0)=g'(0)=0$，设 $f(x)=g(x)/x$（$x\ne0$），$f(0)=0$，证明 $f'(x)$ 在 $x=0$ 处连续。\pt{我的解答}
当$x\ne0$时，$f'(x)=\dfrac{g'(x)x-g(x)}{x^2}$\\ $\displaystyle\lim_{x\to0^-}f'(x)=\lim_{x\to0^+}f'(x)=\lim_{x\to0}\dfrac{g''(x)x+g'(x)-g'(x)}{2x}=\lim_{x\to0}\dfrac{g''(x)}2=\dfrac{g''(0)}2$\\ $g(x)=g(0)+g'(0)x+\dfrac{g''(0)}2x^2+o(x^2)$\\ $\Rightarrow g''(0)=\dfrac2{x^2}\bigl(g(x)-g'(0)x-g(0)+o(x^2)\bigr)=0$。\quad\problem{不能由此推出 $g''(0)=0$，且尚未计算 $f'(0)$。}\\ $\Rightarrow ?$\pt{正确解答}
前两行保留。\fix{由导数定义补上：}$\displaystyle f'(0)=\lim_{x\to0}\frac{f(x)-f(0)}x=\lim_{x\to0}\frac{g(x)}{x^2}=\frac{g''(0)}2$。于是 $\lim_{x\to0}f'(x)=f'(0)$，故 $f'$ 在 $0$ 处连续。
\section*{4.8}\pt{原题}求 $f(x)=x^2\ln(1+x)$ 在 $x=0$ 处的 $n$ 阶导数 $f^{(n)}(0)$（$n\ge3$）。\pt{我的解答}
$\ln(1+x)=\displaystyle\sum_{m=1}^{\infty}(-1)^{m-1}\dfrac{x^m}{m}=\sum_{m=0}^{\infty}(-1)^m\dfrac{x^{m+1}}{m+1}$\\ $\Rightarrow x^2\ln(1+x)=\displaystyle\sum_{m=0}^{\infty}(-1)^m\dfrac{x^{m+2}}{m+1}$。\quad\problem{乘以 $x^2$ 后，$x^{m+1}$ 应为 $x^{m+3}$；这是最早错误。}\\ $=\displaystyle\sum_{n=0}^{\infty}\dfrac{f^{(n)}(0)}{n!}x^n$\\ 令$n=m+2$，则$m=n-2$。\quad\problem{换指标承接了上一行的错误指数；修正后应为 $n=m+3$。}\\ $\Rightarrow f^{(n)}(0)=(-1)^{n-2}\dfrac{n!}{n-1},\ n\ge3$\pt{正确解答}
\fix{只改指数和换指标：}$\displaystyle x^2\ln(1+x)=\sum_{m=0}^{\infty}(-1)^m\frac{x^{m+3}}{m+1}$。令 $n=m+3$，则 $m=n-3$、$m+1=n-2$，所以 $\displaystyle f^{(n)}(0)=(-1)^{n-3}\frac{n!}{n-2}=(-1)^{n-1}\frac{n!}{n-2}$，$n\ge3$。
\end{document}
